\documentclass[lettersize,journal]{IEEEtran}
\usepackage{amsmath,amsfonts,amssymb}
\usepackage{algorithmic}
\usepackage{algorithm}
\usepackage{array}
\usepackage{tabularx}
\usepackage[caption=false,font=normalsize]{subfig}
\usepackage{textcomp}
\usepackage{stfloats}
\usepackage{url}
\usepackage{verbatim}
\usepackage{graphicx}
\usepackage{cite}
\usepackage{hyperref}
\usepackage{bm}
\usepackage{pifont}
\usepackage{amsthm}
\newtheorem{lemma}{Lemma}
\newtheorem{proposition}{Proposition}
\newtheorem{theorem}{Theorem}
\newtheorem{corollary}{Corollary}
\theoremstyle{remark}
\newtheorem{remark}{Remark}
\usepackage{xcolor}
\newcommand{\rev}[1]{\textcolor{blue}{#1}}
\newcommand{\todo}[1]{\textcolor{red}{\textbf{[TODO: #1]}}}
\newcommand{\draftnote}[1]{\par\noindent\textcolor{red}{\textbf{$\blacktriangleright$ DRAFT: #1}}\par\medskip}
\newcommand{\figph}[2]{\fcolorbox{red}{gray!8}{\parbox[c][#1][c]{0.96\linewidth}{\centering
\color{red!65!black}\small\textbf{[figure placeholder]}\\[3pt] #2}}}
\graphicspath{{figs/}}
\hyphenation{op-tical net-works semi-conduc-tor IEEE-Xplore Koop-man}

% ---- method-name macros (freeze notation here) ----
\newcommand{\mbdtrue}{\textsc{MBD-true}}
\newcommand{\dkmbd}{\textsc{DK-MBD}}
\newcommand{\dksplit}{\textsc{DK-MBD-split}}
\newcommand{\bkmbd}{\textsc{BK-MBD}}

\begin{document}

\title{Real-Time Model-Based Diffusion Control\\
with Bilinear Deep Koopman Rollouts}
% \todo{Provisional title -- revisit after abstract is written.}

% Author intentionally omitted for blind submission.
% \author{Author Name(s)\thanks{Affiliation / funding.}}

\maketitle

\begin{abstract}
Model-Based Diffusion (MBD) solves optimal control problems by annealed,
sampling-based score ascent, but it evaluates its rollout model thousands
of times per control step, restricting it to platforms where a physics
simulation is affordable at the control rate. We propose \bkmbd{}, which
replaces the simulator rollout with a learned \emph{bilinear} deep
Koopman surrogate, augmented by a quantile-fitted model-error tube, at a
per-step cost fixed by the lifting dimension. The bilinear form is
essential: we show that whenever a configuration-dependent gain, such
as a vehicle's heading rotation or a manipulator's Jacobian, lies
between the command and the tracked output, the constant input map of a
\emph{linear} Koopman model incurs an irreducible one-step error equal
to the dispersion of that gain. Closed-loop experiments on SE(2) parking and on
end-effector reaching with a kinematic $7$-DOF arm and a physics-simulated
Franka FR3 confirm the analysis in both directions. The linear rollout
fails or stalls ($0/20$, $0/35$, and $7/35$ strict on the FR3), whereas
the bilinear rollout matches the true-dynamics oracle in all $90$ trials
and reduces the FR3 planning step from $4.1$\,s to $64$\,ms, within
$1.3\times$ of the control period. An optimizer ablation attributes the
final-precision gains to the annealed schedule itself.
\end{abstract}

\begin{IEEEkeywords}
Koopman operator, sampling-based control, model-based diffusion, data-driven
control, model learning for control.
\end{IEEEkeywords}

%==============================================================
\section{Introduction}\label{sec:intro}
%==============================================================
\IEEEPARstart{S}{ampling-based} optimal control has become a standard
choice for robotic systems with complex nonlinear dynamics. Instead of solving a
deterministic program at every step, controllers such as Model Predictive
Path Integral (MPPI) control~\cite{kappen2005, theodorou2010,
williams2017mppi, williams2016aggressive} and, more recently, Model-Based
Diffusion (MBD)~\cite{pan2024model, xue2025full} evaluate hundreds of
sampled control sequences by forward simulation and update the nominal
sequence by a cost-weighted average.
MBD reinterprets this population-and-reweight primitive as one step of
reverse diffusion: in this view, the weighted mean is a Monte-Carlo
estimate of the score of a Gaussian-smoothed target
distribution~\cite{song2019generative, song2020score}, and repeating it
along a decreasing noise schedule yields an annealed, coarse-to-fine
search. Because the update requires neither gradients nor structural
assumptions on the dynamics or the cost, and parallelizes naturally on
modern hardware, the approach is broadly applicable. However, the
\emph{rollout model}, evaluated thousands of times at every control
step, carries almost the entire computational burden and thus dictates
the attainable control rate. MBD in particular
inherits its rollouts from a physics simulator~\cite{pan2024model}, which
confines it to offline trajectory optimization or to platforms where a
full simulation step is affordable at the sampling rate.

Koopman operator theory~\cite{koopman1931, mezic2005} provides a
principled way to alleviate this bottleneck by embedding the state in a
space of observables, classically hand-designed~\cite{williams2015edmd,
proctor2016dmdc} and more recently parametrized by deep
networks~\cite{lusch2018deep}, in which the evolution is (near-)linear.
A candidate trajectory can then be propagated by a fixed cascade of
matrix--vector products whose expense is set by the embedding size,
\emph{independently of the complexity of the underlying dynamics}.
\emph{Linear} lifted predictors dominate this literature, since, within
model predictive control (MPC), only a linear predictor yields a convex
program by construction~\cite{korda2018mpc, strasser2026}. MPPI, in
turn, has recently been paired with a learned linear deep Koopman
model, with reported control performance close to true-dynamics MPPI at
a fraction of its cost~\cite{mppidk2026}.

The linearity of the lift, however, forces the input to act through
one constant matrix that does not change with the robot's
configuration. This premise fails for a large class of robots. A
wheeled vehicle's speed command displaces it along its current heading,
and a manipulator's joint-velocity command moves the end-effector
through the Jacobian evaluated at its present posture. In both cases
the effect of the command on the tracked output is modulated by a
\emph{configuration-dependent gain}, and this modulation cannot be
represented by any constant matrix. Notably, in the experimental
settings of the linear-lift controllers above, the model is never
required to represent such a coupling, because the inputs (body-frame
forces or accelerations, with the attitude kinematics integrated
analytically) enter the learned dynamics without state dependence.
We show in Sec.~\ref{sec:method} that the restriction is a model-class
property: the best constant surrogate is the \emph{mean} of the gain
along the visited states, leaving an irreducible one-step error equal
to the gain's dispersion. Sec.~\ref{sec:exp} then shows that its
closed-loop consequence is severe: the linear rollout fails outright on
an analytic SE(2) testbed, stalls short of the strict tolerance on a
physics-simulated manipulator, and recovers only when the coupling is
removed from the learning problem by hand-supplied analytic structure,
which the model-free setting withholds.

For control-affine systems, Koopman theory itself supplies the repair:
the exact Koopman form carries a state-dependent input
map~\cite{iacob2024}, whose smallest realization is the bilinear class.
Bilinear realizations have shown clear modeling and control
advantages~\cite{bruder2021bilinear, folkestad2021, mamakoukas2019},
although these demonstrations lie within MPC, where the bilinear term
destroys the convexity that motivates the linear restriction in the
first place.
A sampling-based optimizer, by contrast, involves no convex program, so
this restriction is no longer necessary: the bilinear model can serve
directly as the rollout, at a matrix--vector cost comparable to that of
the linear model.
Hence, in this paper, we propose \bkmbd{}, a sampling-based control
framework that embeds a learned \emph{bilinear} deep Koopman rollout
inside the annealed diffusion-style update of MBD. At every control
step, the current output is lifted once, every candidate control
sequence is evaluated in the lifted space through the bilinear
recursion, and the annealed update refines the nominal sequence over a
decreasing noise schedule before the first control is executed. Since
the same lifted rollout could equally drive a single-stage MPPI update,
we also ask what the annealed schedule itself contributes, and answer
this question by ablation in Sec.~\ref{subsec:exp_anneal}. A scalar error tube, fitted
to the model's one-step residuals in the spirit of proportional error
bounds for certified Koopman control~\cite{strasser2026, safedmd2024}, is
propagated along each candidate and added to the cost, discouraging the
search from regions where the surrogate is least trustworthy.
Unlike diffusion models used as \emph{learned} motion
priors~\cite{janner2022planning, chi2025diffusion, carvalho2023motion}, no
trajectory distribution is learned here: the diffusion machinery serves
purely as an optimizer, and the learning problem concerns the rollout
dynamics.

The contributions of this paper are as follows:
\begin{itemize}
\item \textbf{\bkmbd{}: MBD at the timescale of the control period.} We
embed a learned bilinear deep Koopman rollout inside the annealed MBD
update, with a quantile-fitted error tube keeping the planner aware of
the surrogate's trust region. The lifted rollout decouples the planning
cost from the complexity of the true dynamics: it is $64\times$ cheaper
than the physics-engine oracle on the FR3, turning a $4.1$\,s planning
step into a $64$\,ms one at no loss in final accuracy.
\item \textbf{Controlled closed-loop validation in both directions.}
Within a single MBD stack that varies only the rollout model, the linear
rollout fails or stalls ($0/20$ and $0/35$ on kinematic testbeds; $7/35$
strict on a physics-simulated Franka FR3), whereas \bkmbd{} matches the
true-dynamics oracle in all $90$ trials at oracle-level final errors
with no analytic model. A structure-informed linear baseline granted
the analytic coupling also succeeds ($90/90$), which isolates the
failure to the model class and delimits what the linear lift would
require in deployment.
\item \textbf{Isolating the contribution of the annealed schedule.}
Holding the learned surrogate fixed and replacing only the update
schedule, we measure what the annealed, multi-stage schedule of MBD
adds over the single-stage MPPI operating point. Strict success climbs
monotonically ($25 \to 28 \to 32 \to 35$ of $35$), and a
Gaussian-homotopy and precision-floor analysis shows that the schedule
matters exactly when the required precision approaches the noise floor
of a fixed-scale sampler (Sec.~\ref{subsec:exp_anneal}).
\item \textbf{Representability analysis of the linear lift.} Supporting
the design, we trace the linear rollout's failure to a model-class
property: the best constant input map is the mean of the
configuration-dependent gain, with an irreducible one-step error equal
to the gain's dispersion. The bilinear lift restores exactly the
structure that the linear lift discards.
\end{itemize}

%==============================================================
\section{Preliminaries}\label{sec:prelim}
%==============================================================

\begin{table}[t]
\centering
\caption{Notation and Symbols}
\label{tab:notation}
\begin{tabularx}{\linewidth}{l X}
\hline
Symbol & Description \\
\hline
$s \in \mathbb{R}^n$, $u \in \mathcal{U} \subset \mathbb{R}^m$ & system state and control input \\
$b = b(s) \in \mathbb{R}^d$ & tracked output / Koopman training state \\
$U = (u_0, \dots, u_{T-1})$ & control sequence over horizon $T$ \\
$J(U)$ & task cost of the rolled-out sequence \\
$z = \Psi_\theta(b) \in \mathbb{R}^r$ & lifted state, $r \ge d$ \\
$C$ & decoder, $Cz = z_{1:d} = b$ \\
$A, B_0, B_i$ & lifted dynamics matrices \\
$M(u) = A + \sum_i u_i B_i$ & input-dependent transition matrix \\
$\sigma_s,\ \alpha$ & annealing noise scale, weight temperature \\
$e_k,\ c_x,\ c_u,\ \beta_e$ & error tube and its constants / weight \\
\hline
\end{tabularx}
\end{table}

\subsection{Problem Formulation}\label{subsec:problem}

Table~\ref{tab:notation} summarizes the notation used throughout the
paper. Consider the discrete-time system
\begin{equation}
s_{t+1} = f(s_t, u_t), \qquad
s_t \in \mathcal{X} \subset \mathbb{R}^{n},\quad
u_t \in \mathcal{U} = [\underline{u}, \bar{u}]^{m},
\label{eq:sys}
\end{equation}
in the \emph{model-free, real-time} setting: the dynamics $f$ cannot be
evaluated at the control rate, whether because they are unidentified or
because their simulation is too slow. The true dynamics therefore serve
only as an offline \emph{oracle} reference and cannot themselves be
deployed as the controller, and model learning must proceed from
interaction data $\mathcal{D} = \{(s_i, u_i, s_i^{+})\}$ alone.

Each task defines observable features $b(s) \in \mathbb{R}^{d}$ used both as
the Koopman training state and in the cost (unicycle:
$b = [x, y, \sin\theta, \cos\theta - 1]$; manipulator:
$b = [q, p_{\mathrm{tcp}}]$, joints and tool-center-point position). The controller operates in a receding
horizon: at state $s_t$ it seeks a control sequence
$U \in \mathbb{R}^{mT}$ minimizing
\begin{equation}
J(U; s_t) = \sum_{k=1}^{T} c\big(\hat{b}_k, u_{k-1}\big)
+ \phi\big(\hat{b}_T\big),
\label{eq:cost}
\end{equation}
where $\hat b_k$ is produced by the rollout model under study, applies the
first control $u_0$ to \eqref{eq:sys}, and replans.

\subsection{Model-Based Diffusion: A Review}\label{subsec:mbd}

MBD~\cite{pan2024model} treats \eqref{eq:cost} as sampling from the
Boltzmann target
\begin{equation}
p_0(U) \;\propto\; \exp\!\big(-J(U)/\alpha\big)
\label{eq:target}
\end{equation}
and descends an annealing ladder of Gaussian-smoothed marginals
\begin{equation}
p_{\sigma}(U) = \big(p_0 * \varphi_{\sigma}\big)(U)
= \mathbb{E}_{\varepsilon \sim \mathcal{N}(0, \sigma^{2} I)}
\big[\, p_0(U - \varepsilon) \,\big],
\label{eq:smooth}
\end{equation}
with $\sigma_1 > \dots > \sigma_S$. By Tweedie's formula, the score of the
smoothed target is a denoising mean shift,
\begin{equation}
\nabla_U \log p_{\sigma}(U)
= \frac{\mathbb{E}\big[\,U_0 \mid U\,\big] - U}{\sigma^{2}} .
\label{eq:tweedie}
\end{equation}
Because the smoothing kernel $\varphi_\sigma(U - U')$ is itself the density
of the proposal $U' \sim \mathcal{N}(U, \sigma^{2} I)$, the posterior mean
in \eqref{eq:tweedie} is a ratio of expectations under the proposal, and
with $N$ candidates $U_k = U + \sigma \varepsilon_k$,
$\varepsilon_k \sim \mathcal{N}(0, I)$, projected onto $\mathcal{U}^{T}$,
the self-normalized estimator is the cost-weighted mean
\begin{equation}
\widehat{\mathbb{E}}\big[U_0 \mid U\big] = \sum_{k=1}^{N} w_k U_k,
\quad
w_k = \frac{e^{-(J(U_k) - J_{\min})/\alpha}}
           {\sum_{j} e^{-(J(U_j) - J_{\min})/\alpha}} .
\label{eq:weights}
\end{equation}
Under the box projection the candidates are no longer exactly Gaussian
about $U$, so \eqref{eq:weights} estimates the score of the
correspondingly projected smoothing rather than \eqref{eq:smooth} itself;
every compared method shares the identical projection.
The update at stage $s$ is a preconditioned score step with optional
Langevin noise, \emph{with the step size tied to the noise level}:
\begin{equation}
U \leftarrow U + \eta_s\, \widehat{\nabla_U \log p_{\sigma_s}}(U)
+ \sqrt{2 \eta_s}\, \xi,
\qquad
\eta_s = \eta\, \sigma_s^{2} .
\label{eq:update}
\end{equation}
With $\eta = 1$ and no Langevin noise, \eqref{eq:update} reduces exactly to
the weighted-mean update $U \leftarrow \sum_k w_k U_k$.

%==============================================================
\section{Proposed Method}\label{sec:method}
%==============================================================

\subsection{Deep Koopman Lifting and Rollout Models}\label{subsec:models}

All learned rollout models share a deep lift with an exact-decoder
convention. With $d = \dim b$,
\begin{equation}
z = \Psi_\theta(b) =
\begin{bmatrix} b \\ \psi_\theta\big(\tau(b)\big) \end{bmatrix}
\in \mathbb{R}^{r},
\qquad
C z = z_{1:d} = b,
\label{eq:lift}
\end{equation}
where $\psi_\theta$ is a small $\tanh$ network and $\tau$ a fixed task
input map (unicycle: $\tau = \mathrm{id}$, $r = 8$; FR3:
$\tau(b) = [\sin q, \cos q]$, $r = 20$). Slicing the decoder avoids
learning a fragile reconstruction and ties the first $d$ lifted coordinates
to the physical output. Two lifted dynamics classes are compared:
\begin{align}
\textbf{linear (DK):}\quad
& z^{+} = A z + B u ,
\label{eq:lin}\\[2pt]
\textbf{bilinear (BK):}\quad
& z^{+} = A z + B_0 u + \sum_{i=1}^{m} u_i B_i z ,
\label{eq:bil}
\end{align}
where, collecting the state-multiplying terms into
$M(u) = A + \sum_i u_i B_i$, one step of \eqref{eq:bil} reads
$z^{+} = M(u)\, z + B_0 u$: the transition matrix acquires an input
dependence, while each step still consists of a small, fixed number of
matrix--vector products. For control-affine systems the exact Koopman
form carries a state-dependent input map~\cite{iacob2024}; the bilinear
class is its smallest realization, restricting that dependence to be
linear in the lifted state. The two classes differ in their input
tangent maps: for \eqref{eq:bil},
$\partial z^{+} / \partial u_i = B_0 e_i + B_i z$ varies with the
lifted state through $B_i z$, whereas the tangent map of the linear
form is the fixed matrix $B$. At planning time the
current output is lifted once, $z_0 = \Psi_\theta(b(s_t))$, and every
candidate is propagated entirely in the lifted space.

\subsection{Why a Constant Input Map Cannot Represent the Coupling}\label{subsec:representability}

Suppose the increment of the tracked output is first-order in the input,
$\Delta y = \Phi(s)\, u\, \Delta t + O(\Delta t^{2})$, where the gain
$\Phi(s)$ varies across configurations (the heading rotation for the
unicycle; the end-effector rows of the Jacobian $J(q)$ for the
manipulator). A linear lift routes the input to the output exclusively through
constant matrices: $CB$ at one step, and the Markov parameters
$C A^{k} B$ at every lag. Since none of these matrices depends on the
state, the full input--output response of \eqref{eq:lin} is linear
time-invariant, however expressive the lifted drift may be.

\begin{proposition}[Irreducible error of a constant gain]
\label{prop:const}
Let inputs satisfy $\mathbb{E}[u u^{\top}] = \lambda I$ independently of
the state $s \sim \rho$. Then
\begin{equation}
\arg\min_{X}\, \mathbb{E}_{s,u}\big\|(\Phi(s) - X) u\big\|^{2}
= \mathbb{E}_{s \sim \rho}\big[\Phi(s)\big],
\label{eq:meangain}
\end{equation}
and the minimum equals
$\lambda\, \mathbb{E}_{s \sim \rho}\big\| \Phi(s) -
\mathbb{E}[\Phi] \big\|_{F}^{2}$.
\end{proposition}
\begin{proof}
$\mathbb{E}_u \|(\Phi - X) u\|^{2}
= \operatorname{tr}\big((\Phi - X)\, \mathbb{E}[u u^{\top}]\,
(\Phi - X)^{\top}\big) = \lambda \|\Phi(s) - X\|_{F}^{2}$; minimizing its
expectation over $s$, a quadratic in $X$, gives the mean, with the variance
as the residual.
\end{proof}

The best constant surrogate is the \emph{mean gain}, and the irreducible
error equals the \emph{dispersion of the gain along the visited
distribution}---a model-class property that no amount of training can
remove. The unicycle is the extreme case: over a parking maneuver the
heading traverses the entire circle, so
$\mathbb{E}[(\cos\theta, \sin\theta)] \approx 0$ and the best constant
input channel for the speed command is \emph{zero}---the
model discards the actuator. The bilinear class instead realizes the
state-dependent gain $u \mapsto C\big(\sum_i u_i B_i z\big)$, whose
columns $C B_i z$ are linear in $z$. The heading rotation is therefore
exactly expressible, since $\sin\theta$ and $\cos\theta$ are (affine)
coordinates of $z$, and the manipulator Jacobian, whose entries are
multilinear in $\{\sin q_i, \cos q_i\}$ across joints, is expressible
once the learned features supply the cross-products. We emphasize that
Proposition~\ref{prop:const} is a representability statement about the
one-step, first-order input channel under the stated excitation, not a
closed-loop theorem. How this representability gap is amplified into
closed-loop failure once an optimizer selects controls through the
model is traced empirically in Sec.~\ref{subsec:exp_mechanism}.

\subsection{The Structure-Informed Alternative and Its Cost}\label{subsec:split}

Proposition~\ref{prop:const} indicates that the only repair available
to the linear class is to remove the coupling from the learning
problem. We formalize
this as the \dksplit{} baseline, which mirrors the experimental
structure of linear-lift controllers~\cite{mppidk2026, korda2018mpc}.
For the unicycle, a linear map $G$ is fitted by least squares from the
input to the \emph{body-frame} displacement,
$[\Delta p_{\mathrm{body}}; \Delta\theta] = G u$ with
$\Delta p_{\mathrm{body}} = R(-\theta_k)(p_{k+1} - p_k)$, and rollouts
compose the prediction analytically with the current heading (an
$SE(2)$ composition). For the manipulator, a linear deep Koopman model
is trained on the joint state only, and the analytic forward kinematics
is applied to the decoded joints when evaluating the cost. In both cases the input then
enters the \emph{learned} part state-independently, and
Proposition~\ref{prop:const} no longer binds. \dksplit{} is the fair
representative of what linear-lift control actually deploys, and its
requirement is precisely the analytic rotation or forward kinematics
that the model-free setting withholds. The bilinear rollout
\eqref{eq:bil} needs neither.

\subsection{Multi-Step Training}\label{subsec:training}

The planner consumes $T$-step rollouts, so the surrogate must remain
accurate over a horizon rather than at a single transition. Models
fitted by plain one-step regression drift away from the data
distribution before the planning horizon is reached. All learned models are
therefore trained on random-rollout snippets by the same $H$-step loss
($H = 15$),
\begin{equation}
\mathcal{L} = \frac{1}{H} \sum_{h=1}^{H}
\Big[ \big\| C \hat{z}_h - b_h \big\|^{2}
+ \gamma \big\| \hat{z}_h - \Psi_\theta(b_h)^{\mathrm{sg}} \big\|^{2}
\Big],
\label{eq:loss}
\end{equation}
with $\hat z_0 = \Psi_\theta(b_0)$, $\hat z_{h+1}$ given by
\eqref{eq:lin} or \eqref{eq:bil}, $\gamma = 0.1$, and stop-gradient
($\mathrm{sg}$) on the latent target. On the 7-DOF manipulator the $m$
per-channel matrices make the bilinear recursion prone to instability
during training. Setting $B_i = 0$ at initialization resolves this
issue, because the optimization then begins at the linear model and
introduces bilinear structure only where the data requires it.
Datasets are generated once and shared across every learned model and
training seed, so that the compared classes differ in nothing but the
lifted dynamics \eqref{eq:lin} vs.\ \eqref{eq:bil}.

\subsection{Model-Error Tube}\label{subsec:tube}

To make the planner aware of where the surrogate can be trusted, the
one-step lifted residuals on the training set,
$r_k = \Psi_\theta(b_{k+1}) - \big(M(u_k)\Psi_\theta(b_k) + B_0 u_k\big)$,
are covered by a proportional bound
$\|r_k\| \le c_x \|z_k\| + c_u \|u_k\|$ fitted at quantile coverage. A
least-squares fit of $(\hat c_x, \hat c_u)$ to the residual norms is
rescaled by the empirical $(1-\delta)$-quantile of the coverage ratios,
\begin{equation}
\begin{aligned}
s_\delta &= Q_{1-\delta}\!\left[
\frac{\|r_k\|}{\hat c_x \|z_k\| + \hat c_u \|u_k\|}
\right], \\
(c_x, c_u) &= \max(s_\delta, 1) \,(\hat c_x, \hat c_u),
\end{aligned}
\label{eq:tubefit}
\end{equation}
with $\delta = 10^{-3}$. Quantile coverage, rather than a worst-case fit
over every residual, keeps the constants proportional to the typical
model error. Sec.~\ref{subsec:exp_cost} quantifies how a worst-case fit
lets a few rest-to-motion transients dominate the planning cost.
During planning a scalar tube is propagated along each candidate with the
precomputed norm bound
$\bar m(u) = \|A\|_2 + \sum_i |u_i| \|B_i\|_2 \ge \|M(u)\|_2$:
\begin{equation}
e_{k+1} = \big(\bar m(u_k) + c_x\big) e_k
+ c_x \|\hat z_k\| + c_u \|u_k\|,
\qquad e_0 = 0 .
\label{eq:tube}
\end{equation}
On any segment where the fitted bound holds, an induction over the exact
error dynamics $\delta_{k+1} = M(u_k)\delta_k + r_k$ of the bilinear
rollout shows that the recursion dominates the true-vs-predicted lifted
error: $e_k \ge \|\delta_k\|$ with
$\delta_k = \Psi_\theta(b_k) - \hat z_k$.
The planner then minimizes the tube-augmented cost
\begin{equation}
\tilde J(U) = J(U) + \beta_e \sum_{k} e_k ,
\label{eq:tubecost}
\end{equation}
where $\beta_e$ scales with the task cost magnitude; switching the
penalty off ($\beta_e = 0$) leaves the plain \bkmbd{} planner. The
constants $(c_x, c_u)$ are coverage estimates fitted from data, and we
attach no identification-theoretic certificate to them.

In summary, a \bkmbd{} control step proceeds as follows. The measured
output is encoded once via \eqref{eq:lift}. Each annealing stage
$s = 1, \dots, S$ then draws $N$ candidate sequences, scores them by
running \eqref{eq:bil} in the lifted space while accumulating
\eqref{eq:tube}, forms the weights \eqref{eq:weights} on the
tube-augmented cost \eqref{eq:tubecost}, and applies \eqref{eq:update}
to the nominal sequence. Finally the first control is executed, and the
sequence is shifted by one step to warm-start the next replan.

%==============================================================
\section{Experiments}\label{sec:exp}
%==============================================================
We evaluate the framework on two main testbeds, both first-order
coupled in the sense of Sec.~\ref{subsec:representability}: an SE(2)
unicycle parking task, where the coupling is exact, and end-effector
reaching with a Franka Research 3 (FR3) manipulator simulated with full
rigid-body physics in MuJoCo. On the FR3, the joint-velocity servo
realizes the first-order interface only up to lag and load.
A third testbed, a kinematic $7$-DOF arm, carries the optimizer study of
Sec.~\ref{subsec:exp_anneal}. The
experiments are organized to support the claims of Sec.~\ref{sec:intro}:
the failure of the full-lift linear rollout and the oracle-level
performance of the bilinear rollout
(Secs.~\ref{subsec:exp_unicycle}--\ref{subsec:exp_franka}), the mechanism
by which a modest open-loop gap becomes that closed-loop failure once an
optimizer sits on top of the model (Sec.~\ref{subsec:exp_mechanism}),
the contribution of the annealed schedule over a single-stage sampler
(Sec.~\ref{subsec:exp_anneal}), the
rollout-cost decoupling together with an empirical validation of the tube
(Sec.~\ref{subsec:exp_cost}), and finally the boundary of the regime in
which the bilinear class is needed at all
(Sec.~\ref{subsec:exp_regime}).

\subsection{Protocol and Compared Methods}\label{subsec:setup}

All four planners run inside one implementation of the MBD optimizer of
Sec.~\ref{subsec:mbd}. The \emph{only} component that differs across
methods is the rollout model that maps a candidate control sequence to
a cost. \mbdtrue{} rolls out the true dynamics (the analytic model on the
unicycle, the MuJoCo simulator on the FR3) and serves as the oracle
upper bound. \dkmbd{} rolls out the full-lift linear deep Koopman model
\eqref{eq:lin}. \dksplit{} is the structure-informed control of
Sec.~\ref{subsec:split}: it is granted the analytic coupling (the SE(2)
group action on the unicycle, the forward kinematics on the FR3) and
learns only the input-affine remainder. \bkmbd{} rolls out the bilinear
model \eqref{eq:bil} with the tube cost \eqref{eq:tubecost}. The optimizer
uses $S = 5$ annealing stages, temperature $\alpha = 0.4$, and step scale
$\eta = 1$ (the deterministic weighted-mean limit of \eqref{eq:update});
per task we fix the sample count, noise schedule, horizon, and cost across
all methods (unicycle: $N = 768$, $\sigma\!: 0.9 \to 0.2$, $T = 40$;
FR3: $N = 800$, $\sigma\!: 1.2 \to 0.3$, $T = 15$).

The learned models follow one training recipe
(Sec.~\ref{subsec:training}). For each of five seeds, the linear,
bilinear, and split models are trained on the same dataset of $6{,}000$
random-excitation snippets of length $15$, so that every learned-model
comparison shares data, architecture, epochs, and optimizer, and
differs only in the lifted dynamics class. Every
(task, seed) pairing then runs all closed-loop cases from identical initial
sequences, giving $5 \times 4 = 20$ trials on the unicycle and
$5 \times 7 = 35$ trials on the FR3 (and likewise $35$ on the kinematic
arm of Sec.~\ref{subsec:exp_anneal}). We report the task success rate, the
strict success rate at a tighter tolerance, the final error, and the
planning time per control step.

\subsection{Unicycle Parking}\label{subsec:exp_unicycle}

The unicycle state is $(x, y, \theta)$ with speed and turn-rate inputs,
integrated by RK4 at $\Delta t = 50$\,ms. The task is to park at the
origin from four start poses within $180$ steps, with success defined
as a final position error below $0.3$\,m (strict: $0.08$\,m position
\emph{and} $0.2$\,rad heading). The base features used for lifting and cost are
$[x, y, \sin\theta, \cos\theta - 1]$, and the coupling
$R(\theta)\,[v, 0]^\top$ is exactly the configuration-dependent gain of
Proposition~\ref{prop:const}.

\begin{table}[!t]
\centering
\caption{Unicycle parking, 5 seeds $\times$ 4 cases.}
\label{tab:unicycle}
\begin{tabular}{lcccc}
\hline
Method & Park & Strict & Final err.\ [m] & ms/step \\
\hline
\mbdtrue{} & $20/20$ & $20/20$ & $0.070 \pm 0.008$ & $81$ \\
\dkmbd{} & $0/20$ & $0/20$ & $2.383 \pm 1.706$ & $43$ \\
\dksplit{} & $20/20$ & $20/20$ & $0.066 \pm 0.010$ & $58$ \\
\bkmbd{} (ours) & $20/20$ & $18/20$ & $0.071 \pm 0.026$ & $95$ \\
\hline
\end{tabular}
\end{table}

\begin{figure}[!t]
\centering
\includegraphics[width=\linewidth]{figs/unicycle_trajectories.png}
\caption{Unicycle parking trajectories for all methods over five seeds
(gray disk: $0.3$\,m park radius). \dkmbd{} (blue) veers away in every
trial, while \dksplit{} (orange) and \bkmbd{} (red) are indistinguishable
from the oracle (green).}
\label{fig:unicycle_traj}
\end{figure}

\dkmbd{} does not merely park imprecisely but diverges, ending with a
mean final error of $2.38$\,m on a task whose starts are $2$--$3$\,m
from the goal. Its constant input map commits to one average heading
direction (Proposition~\ref{prop:const}), and the annealed optimizer
then faithfully optimizes the wrong model. Both control
conditions behave as predicted (Table~\ref{tab:unicycle},
Fig.~\ref{fig:unicycle_traj}): granting the linear model the SE(2)
structure (\dksplit{})
or replacing the model class by the bilinear lift (\bkmbd{}) restores
oracle-level parking, with final errors within one standard deviation of
\mbdtrue{}.

\subsection{Franka FR3 Reaching in Physics Simulation}\label{subsec:exp_franka}

The second testbed is a Franka
Research 3 manipulator with attached Franka Hand, simulated in MuJoCo with
full rigid-body dynamics. The robot is controlled through a
gravity-compensated joint-velocity servo interface, the same
abstraction exposed by the vendor controller on the physical robot, at
a control period of $50$\,ms spanning $25$ physics substeps. The
surrogate must therefore absorb servo lag and the inertial load of the
hand rather than an exact integrator. The control point is the tool center point (TCP) between the
two fingertips, and the Koopman state is $b = [q, p_{\mathrm{tcp}}]$
with measured joint positions. At every control step the planner
searches a $105$-dimensional decision space ($m = 7$ joint velocities
over $T = 15$ steps). Training data are $6{,}000$ coherent random
joint-velocity snippets rolled through the simulator; seven reach targets
are attempted for up to $120$ steps ($6$\,s), with success below $5$\,cm
and strict success below $2.5$\,cm. The oracle plans through the simulator
itself with threaded batch rollouts.

\begin{table}[t]
\centering
\caption{FR3 reaching in MuJoCo, 5 seeds $\times$ 7 targets.}
\label{tab:franka}
\begin{tabular}{lcccc}
\hline
Method & Reach & Strict & Final err.\ [m] & ms/step \\
\hline
\mbdtrue{} & $35/35$ & $35/35$ & $0.023 \pm 0.002$ & $4094$ \\
\dkmbd{} & $13/35$ & $7/35$ & $0.054 \pm 0.022$ & $37$ \\
\dksplit{} & $35/35$ & $35/35$ & $0.023 \pm 0.001$ & $56$ \\
\bkmbd{} (ours) & $35/35$ & $35/35$ & $0.023 \pm 0.002$ & $64$ \\
\hline
\end{tabular}
\end{table}

Table~\ref{tab:franka} shows that the
pattern transfers to full physics. \bkmbd{} again matches the oracle
exactly ($0.023 \pm 0.002$ vs.\ $0.023 \pm 0.002$\,m, all trials
strict) and traces the same direct TCP paths. \dkmbd{}, in contrast, is
unreliable: it wanders across the workspace, reaching the strict
tolerance in only $7/35$
trials and exceeding the loose one in $22/35$. The planning-time column
carries the practical significance of these results. Rolling out the
physics engine costs the oracle $4.1$\,s per $50$\,ms control step, two
orders of magnitude away from real time even with threaded rollouts.
\bkmbd{} plans in $64$\,ms, a $64\times$ reduction that brings the same
annealed optimizer to within $1.3\times$ of the control rate while
losing nothing in final accuracy.

Two alternative explanations for the linear rollout's failure remain to
be examined: insufficient capacity, and the possibility that a generic
nonlinear dynamics model would suffice. Table~\ref{tab:baselines}
addresses both. Regarding capacity, the enlarged linear lift at
$r = 57$ spends more dynamics parameters, and a higher per-step FLOP
count, than the bilinear model at $r = 20$, yet its held-out TCP error
does not move to three decimals; the limitation is the class-level
ceiling of Proposition~\ref{prop:const}, not a parameter shortage.
Regarding the second explanation, an MLP dynamics model given the same
features, multi-step loss, and budget does capture the coupling, and
its open-loop accuracy even slightly exceeds the bilinear model's. In
closed loop, however, it stalls at the step limit $2.9$--$4.7$\,cm from
the target and reaches the strict tolerance in only $26/35$ trials
against $35/35$ for the bilinear rollout---a second instance of
open-loop accuracy failing to certify closed-loop precision
(Sec.~\ref{subsec:exp_mechanism}). It also incurs a $2.7\times$ higher
rollout cost and sacrifices the lifted-space linearity that the tube
recursion~\eqref{eq:tube} exploits.
Among these candidates, only the bilinear lift combines
nonlinear-class accuracy, matrix--vector rollout cost, tube-compatible
structure, and oracle-level closed-loop precision.

\begin{table}[t]
\centering
\caption{Ruling out capacity and black-box explanations on the FR3:
held-out open-loop TCP error (training seed $0$), rollout cost of an
$800$-candidate batch, and closed-loop strict success under the
identical planner ($5$ seeds $\times$ $7$ targets; linear row from
Table~\ref{tab:franka}).}
\label{tab:baselines}
\begin{tabular}{lcccc}
\hline
Model & $r$ & $H{=}15$ [m] & ms & Strict \\
\hline
Linear & $20$ & $0.336$ & $0.3$ & $7/35$ \\
Linear (large) & $57$ & $0.336$ & $0.6$ & --- \\
MLP dynamics & --- & $0.236$ & $3.0$ & $26/35$ \\
Bilinear (ours) & $20$ & $0.252$ & $1.1$ & $\mathbf{35/35}$ \\
\hline
\end{tabular}
\end{table}

\begin{figure*}[!t]
\centering
\includegraphics[width=0.9\textwidth]{figs/franka_final_pose_T7_seed0.png}
\caption{Final FR3 configuration (bold kinematic chain) and TCP path
(dotted) on the seventh reach target for each method. \bkmbd{}, the
oracle, and \dksplit{} bring the tool to the target ($\star$) in
$19$--$26$ steps; \dkmbd{} leaves the arm short of the goal at the step
limit ($6.9$\,cm).}
\label{fig:franka_pose}
\end{figure*}

\begin{figure}[!t]
\centering
\includegraphics[width=0.494\linewidth]{figs/bk_mbd_franka.png}\hfill
\includegraphics[width=0.494\linewidth]{figs/dk_mbd_franka.png}
\caption{Executed closed-loop FR3 motions on a representative reach,
rendered as afterimages (translucent snapshots along the trajectory).
\bkmbd{} (left) sweeps directly to the target (green sphere; TCP trace in
red), whereas \dkmbd{} (right) wanders across the workspace without
settling on the target (blue trace).}
\label{fig:franka_ghost}
\end{figure}

\subsection{From Open-Loop Error to Closed-Loop Failure}
\label{subsec:exp_mechanism}

The closed-loop results above raise a question about the mechanism:
does the linear lift fail because it is a much worse \emph{predictor},
or because of how an \emph{optimizer} uses it? Because MBD interrogates the rollout model
$S \times N$ times per control step and executes whatever the model scores
best, the distinction matters, and it can be read off the open-loop error
of the very checkpoints used in closed loop.
Figure~\ref{fig:pred_horizon} shows the open-loop prediction error of the
FR3 checkpoints against the rollout horizon on held-out data, averaged
over the five training seeds and decomposed into the output rows that
carry the configuration-dependent coupling (TCP) and the joint rows,
where the input enters state-independently.

\begin{figure}[t]
\centering
\includegraphics[width=\linewidth]{figs/pred_error_horizon.png}
\caption{Open-loop prediction error vs.\ horizon for the trained FR3
checkpoints used in closed loop (mean $\pm$ std over five seeds,
held-out data, log scale). The classes coincide on the joint rows
(dashed); on the TCP rows (solid) the linear model is only $1.33\times$
worse at the planning horizon $H = 15$---yet reaches the strict
tolerance in only $7/35$ closed-loop trials against $35/35$ for the
bilinear model.}
\label{fig:pred_horizon}
\end{figure}

Two observations localize the failure. First, the penalty is exactly where
Proposition~\ref{prop:const} predicts: on the joint rows the two classes
are indistinguishable, while on the coupled rows the linear model carries
a persistent excess error. Second, and more consequential, the open-loop
gap is \emph{modest}---$1.33\times$ at the planning horizon---because
the Jacobian along reaching motions has a non-zero mean, which the
constant input map of Proposition~\ref{prop:const} captures. The
irreducible part is only the \emph{dispersion} around that mean. In
closed loop, this modest gap becomes $7/35$ strict versus $35/35$. The
amplification is a property of optimizing through the model rather than
predicting with it. The residual left by the mean gain is a
\emph{structural bias} that always points the same way for a given
configuration; the annealed search selects the candidates whose
predicted progress exploits this bias, and receding-horizon replanning
re-commits to the biased direction at every step.
Fig.~\ref{fig:franka_pose} and Fig.~\ref{fig:franka_ghost} visualize
this end state: \dkmbd{} leaves the arm short of the goal when the
step limit expires, while every other method brings the tool to the
target. The unicycle
lies at the opposite end of the same spectrum: a parking maneuver drives the
heading around the entire circle, leaving a near-zero mean gain, so the
one-step open-loop gap is already $40\times$ ($0.082$ vs.\ $0.002$
RMSE) and the closed loop diverges outright. Open-loop error
therefore under-reports the model-class deficit in proportion to how
benign the visited gain distribution is. A rollout surrogate for a
sampling-based optimizer must consequently be judged in closed loop, on
the coupled output rows, rather than by aggregate prediction scores.

\subsection{Contribution of the Annealed Schedule}\label{subsec:exp_anneal}

\begin{table}[!t]
\centering
\caption{Optimizer ablation on the kinematic arm: identical bilinear
checkpoints, cost, and tube; only the update schedule varies ($5$ seeds
$\times$ $7$ targets). (a) is the MPPI operating point; (d) is the full
anneal used by \bkmbd{}.}
\label{tab:anneal}
\begin{tabular}{lcccc}
\hline
Schedule & Reach & Strict & Final err.\ [m] & ms \\
\hline
(a) $1{\times}800$, $\sigma{\equiv}1.2$ & $32/35$ & $25/35$ & $0.029 \pm 0.013$ & $18$ \\
(b) $1{\times}4000$, $\sigma{\equiv}1.2$ & $32/35$ & $28/35$ & $0.028 \pm 0.008$ & $24$ \\
(c) $5{\times}800$, $\sigma{\equiv}1.2$ & $34/35$ & $32/35$ & $0.027 \pm 0.017$ & $71$ \\
(d) $5{\times}800$, $\sigma{:}\,1.2{\to}0.3$ & $\mathbf{35/35}$ & $\mathbf{35/35}$ & $\mathbf{0.024 {\pm} 0.001}$ & $69$ \\
\hline
\end{tabular}
\end{table}

The optimizer itself is the remaining free variable: does the annealed,
multi-stage schedule of MBD contribute anything that a single-stage
sampler---the MPPI operating point---would not achieve with the same
surrogate?
We answer on a third testbed on which the question is sharpest: a
kinematic $7$-DOF arm reaching end-effector targets under
joint-velocity inputs, with the same training recipe and protocol
($5$ seeds $\times$ $7$ targets, success below $5$\,cm, strict below
$2.5$\,cm). Its main comparison mirrors the FR3 pattern: the linear
rollout fails outright ($0/35$, final error $0.73 \pm 0.40$\,m), while
\bkmbd{} matches the oracle ($35/35$, $0.024$ vs.\ $0.023$\,m).
Holding the trained bilinear checkpoints, task, cost, and tube fixed, we
then vary \emph{only} the update schedule (Table~\ref{tab:anneal}).

Strict success climbs monotonically as optimizer structure is added
($25 \to 28 \to 32 \to 35$), but only the full anneal reaches $35/35$,
at the same rollout budget as the fixed-noise five-stage variant (c).
Five times more samples in a single stage (b) yield almost no
improvement.
Two mechanisms explain the ordering.
First, with a learned surrogate the evaluated cost is
$\hat J = J + e$, where the model-error term $e$ inherits fine-scale
structure from error accumulation along rollouts. Each annealing stage
descends a Gaussian-smoothed landscape in which the Fourier components
of $e$ are attenuated by $e^{-\sigma^{2}\|\omega\|^{2}/2}$: early
(large-$\sigma$) stages select the basin on a landscape whose spurious
model-induced minima are suppressed, while late stages refine where
local bias is small. A single fixed scale must serve both roles at
once.
Second, a fixed scale has a \emph{precision floor}: the weighted-mean
estimator carries Monte-Carlo noise of order
$\sigma/\sqrt{N_{\mathrm{eff}}}$ per axis
($N_{\mathrm{eff}} = 1/\sum_k w_k^{2}$). In receding horizon the
control is executed after a single update, so this noise enters the
plant directly.
The failure signature confirms the mechanism: seven of the eight
fixed-noise failures across (a)--(c) are marginal stalls at
$0.050$--$0.079$\,m against the $0.05$\,m threshold---never
divergence---and annealing to $\sigma_S = 0.3$ removes every one.
On the FR3 the same sweep leaves all four configurations at $35/35$
strict ($17$--$64$\,ms per step), since its $2.5$\,cm tolerance lies
above the floor there. The schedule therefore matters exactly when the
required precision approaches the floor, which is the regime a
surrogate-based planner enters whenever tolerances tighten.

\subsection{Rollout Cost and the Error Tube}\label{subsec:exp_cost}

The planning-time columns quantify the cost decoupling claimed in
Sec.~\ref{sec:intro}. A sampling planner spends
$S \times N \times T$ model evaluations per control step. For \bkmbd{}
each evaluation is one bilinear matrix--vector recursion in
$\mathbb{R}^{r}$, so its per-step cost is set by the lifting dimension
alone and is essentially constant across our testbeds ($64$--$95$\,ms).
The oracle's cost instead tracks the complexity of the true dynamics,
from $81$\,ms on the analytic unicycle to $4.1$\,s on the physics
engine---two orders of magnitude the moment rollouts require a
simulator.

The tube cost \eqref{eq:tubecost} acts as a soft trust region during the
anneal, and Fig.~\ref{fig:tube_val} checks empirically that the
quantile-fitted recursion still does the job it is asked to do. On
$2{,}000$ held-out FR3 rollouts of $40$ steps ($80{,}000$ propagated
steps, $2.7\times$ the planning horizon), the tube bounds the true
lifted error in all but $21$ steps ($0.026\%$): the one-step
$1-\delta = 99.9\%$ coverage of \eqref{eq:tubefit} transfers
essentially unchanged through the propagation \eqref{eq:tube}. The
bound is loose
(median $e_k / \|\delta_k\| \approx 1.8 \times 10^{3}$ at the planning
horizon), which is inherent to the norm recursion: with
$\|A\|_2 \approx 1.00$ and $\sum_i \|B_i\|_2 \approx 1.24$ the growth
factor $\bar m(u) + c_x$ exceeds one at every step. For \bkmbd{} this
looseness is immaterial, because the tube enters as a \emph{relative}
penalty that orders candidates by how far they leave the trusted
region, rather than as a constraint. For the same reason, however, we
regard the tube as empirical and attach no certificate semantics to
it.

\begin{figure}[t]
\centering
\includegraphics[width=0.75\linewidth]{figs/tube_validation.png}
\caption{Empirical tube $e_k$ vs.\ the true lifted error on a
representative held-out FR3 rollout, with the quantile-fitted constants
used by the planner. Over $2{,}000 \times 40$ held-out steps the tube is
violated in $0.026\%$ of steps, matching the fitted $99.9\%$ one-step
coverage.}
\label{fig:tube_val}
\end{figure}

The fitting of the tube, however, must respect the residual
distribution. Fitting $(c_x, c_u)$ to the worst-case residual instead
of the quantile bound \eqref{eq:tubefit} lets a handful of
rest-to-motion transients, which a velocity-free lifted state cannot
represent, inflate $c_x$ by $40\times$ (from $0.008$ to $0.32$) on the
FR3. The propagated bound \eqref{eq:tube} then grows exponentially over
the horizon, and the tube term---not the task cost---dominates the
anneal. With no change to the model, \bkmbd{} is left oscillating
$23$\,cm from a target that it otherwise reaches with $2.4$\,cm
accuracy. The quantile fit ($\delta = 10^{-3}$) restores the intended
regularizing role.

\subsection{Delimiting the Regime}\label{subsec:exp_regime}

The bilinear rollout was motivated by first-order coupling. Two compact
studies under the same training recipe locate the boundary of that
motivation, in decreasing order of coupling strength.
The first study considers a $3$D quadrotor in MuJoCo, the canonical
platform of the linear-lift control literature, in which the coupling
acts one integration below the tracked output: the thrust command must
first tilt into an acceleration through the attitude before the
position responds. At a single step the bilinear model is far more
accurate even on the linear-velocity rows, since the one-step increment
$\Delta v \propto T\,R e_3$ \emph{is} a state--input product, and it
remains $2.3$--$2.4\times$ more accurate on the full feature vector at
$H = 12$. Over the horizon, however, both classes are limited by the
same attitude-prediction error, and the velocity-row gap settles at
$1.5$--$1.8\times$ (three training seeds). The remaining gap reflects
\emph{shared} prediction error rather than the directional
input-channel bias of Sec.~\ref{subsec:exp_mechanism}, so a coupling
hidden behind an integrator leaves the linear lift serviceable for
position control. This is the very regime in which the Koopman control
literature reported its linear-lift successes (Sec.~\ref{sec:intro}).
The second study considers a torque-limited pendulum, whose input
involves no configuration-dependent gain, so that
Proposition~\ref{prop:const} imposes no structural penalty: inside the
same annealed planner, the linear class with a $128$-neuron encoder
swings up from hanging on $5/5$ planning seeds. The bilinear class
contributes parameter economy here rather than feasibility. At the same
lifted dimension $r = 8$, it matches with a $32$-neuron encoder the
multi-step accuracy that the linear model needs $256$ neurons to reach
(mean over three training seeds), and this difference carries into
closed loop at tight budgets, where the width-$64$ linear model fails
swing-up ($0/5$) while the bilinear one already succeeds ($4/5$).
Together with
Secs.~\ref{subsec:exp_unicycle}--\ref{subsec:exp_franka}, these studies
delimit the claim of this paper: the bilinear rollout is required when
the tracked output is first-order coupled, optional when the coupling
acts behind an integrator, and an efficiency improvement when the
coupling is absent.

%==============================================================
\section{Conclusion}\label{sec:conclusion}
%==============================================================
This paper presented \bkmbd{}, which replaces the simulator rollouts of
Model-Based Diffusion with a learned bilinear deep Koopman surrogate and
a quantile-fitted model-error tube. The motivating representability
argument states that the best constant input map is the mean of the
configuration-dependent gain, leaving an irreducible error equal to its
dispersion. Closed-loop experiments on an SE(2) unicycle, a kinematic
$7$-DOF arm, and a physics-simulated Franka FR3 confirmed this argument
inside a single planning stack. The full-lift linear rollout failed or
stalled ($0/20$, $0/35$, and $7/35$ strict), while the bilinear rollout
matched the true-dynamics oracle in all $90$ trials and reduced the FR3
planning step from $4.1$\,s to $64$\,ms, within $1.3\times$ of the
$50$\,ms control period. An optimizer ablation further showed that the
annealed schedule itself supplies the final precision that a
single-stage sampler cannot reach.

Two limitations qualify these results. First, the tube constants are
coverage estimates fitted from data, since certified bounds in the
spirit of~\cite{safedmd2024} would require identification machinery
that our deep lift does not yet provide. Second, all our tasks, while
first-order coupled, are free of contact. Certified tube constants and
contact-rich dynamics therefore remain open problems. Conversely, the
linear lift remains a serviceable controller under the same planner
wherever the input acts state-independently or only through an
integrator (Sec.~\ref{subsec:exp_regime}), so the bilinear class is
the repair demanded specifically by first-order coupled outputs.
Deployment on the physical FR3 through the same joint-velocity
interface is planned as the next step.

\bibliographystyle{IEEEtran}
\bibliography{IEEEabrv,references,ours}

\end{document}
